\documentclass{report}
\usepackage[utf8]{inputenc}

\title{Una temprana historia de la Internet}
\author{Arteaga Hernández Alejandro Iabin 313033414}
\date{Septiembre 2018}

\usepackage{natbib}
\usepackage{graphicx}

\begin{document}

\maketitle

Desde tiempos anteriores, siempre se había pensando 
en que tarde o temprano habría, un gran banco de datos mundial, un gran cerebro, una gran red. \\ 
Se empezaba a visualizar la necesidad de conectar computadoras de manera eficiente, y aqui es cuando se empiezan a desarrollar las herramientas que más tarde se unirían para crear la ARPANET\\

Leonard Kleinrock, el autor del artículo , se da cuenta al escribir su tesis que las computadoras pueden ser conectadas y que la manera de canales de comunicación, no era eficiente, ya que se desperdiciaba mucho del canal de comunicación, ya que las conexiones no estaban trabajando todo el tiempo.\\ En 1962 en su tesis doctoral, define la teoría de lo hoy conocido como conmutación de paquetes, diseñado para transmitir mensajes entre computadoras, fragmentándolos en paquetes, y así desarrollo toda la teoría necesaria.\\
Paul Baran por su parte desarrolló una teoría equivalente, aunque se percató que las conexiones telefónicas no eran suficiente, y creo un trabajo donde hablaba de paquetes. \\

Usando el trabajo escrito por Leonard, Donald Davies buscando desarrollar una nueva red, pidió apoyo al gobierno británico, pero este se negó.\\
Ya con bases de rechazo J. C. R. Licklider había publicado un trabajo llamado ”Man-Computer Symbiosis” que buscaba la interacción mas cercana entre hombre y máquina, así que buscó la forma de implementarlo, Licklider fue nombrado el director de (IPTO) de ARPA, así que con esto Licklider buscó la creación de la red, con los recursos se unen los hilos y se da inicio a la ARPANET.
 \\
 El principal problema era que intentaban conectar hardware incompatible, así que había que hacer protocolos genéricos para la conexión, además de crear los nodos y la conexión física, en 1967 es presentada la ARPANET.\\
 Con el uso de un IMP (Interface Message Processor) que se encargaba de la traducción a paquetes y manejo de la red, en 1969 se conecta el primer IMP el primer nodo, posteriormente se conectaron 4 nodos más.
 
 \includegraphics[scale=0.5]{Capture.PNG}
 Los 4 nodos originales 1969\\
 \\En 1971 con 15 nodos, se desarrolla la aplicación de correo electrónico, en 1972, la ARPANET fue presentada públicamente, después se conectarían mas redes a esta red, y sería conocida como
 Internet.


\end{document}
